\documentclass[12pt,a4paper]{article}

\usepackage[T2A]{fontenc}
\usepackage[utf8]{inputenc}
\usepackage[russian]{babel}
\usepackage{amsmath,amssymb,amsthm}
\usepackage{mathtools}
\usepackage{enumitem}
\usepackage{listings}
\usepackage{xcolor}
\usepackage{array}
\usepackage{hyperref}
\usepackage{stmaryrd}

\lstdefinestyle{code}{
  basicstyle=\ttfamily\small,
  keywordstyle=\color{blue},
  commentstyle=\color{gray},
  stringstyle=\color{red!60!black},
  columns=fullflexible,
  keepspaces=true,
  frame=single,
  breaklines=true,
  showstringspaces=false
}

\lstset{style=code}

\newtheorem{definition}{Определение}
\newtheorem{theorem}{Теорема}
\newtheorem{lemma}{Лемма}
\newtheorem{example}{Пример}
\newtheorem{remark}{Замечание}

\title{Билет №32 \\ \small Процедурные языки программирования. Основные управляющие конструкции, структура программы. Работа с данными: переменные и константы, типы данных (булевский, целочисленные, плавающие, символьные, типы диапазона и перечисления, указатели), структуры данных (массивы и записи). Процедуры (функции): вызов процедур, передача параметров (по ссылке, по значению, по результату), локализация переменных, побочные эффекты. Обработка исключительных ситуаций. Библиотеки процедур и их использование. (СпБГУ)}
\author{Сериков Владислав Александрович}
\date{13 мая 2026}

\begin{document}

\maketitle
\tableofcontents
\newpage

\section{Процедурные языки программирования}

\subsection{Понятие процедурного программирования}

\textbf{Процедурное программирование} --- парадигма программирования, в которой программа рассматривается как последовательность команд, изменяющих состояние вычислительной системы, а основным средством структурирования являются \textbf{процедуры} и \textbf{функции}.

Процедурный язык программирования обычно предоставляет:
\begin{itemize}
    \item переменные и операции присваивания;
    \item встроенные и пользовательские типы данных;
    \item управляющие конструкции: последовательность, ветвление, циклы;
    \item процедуры и функции;
    \item области видимости и правила времени жизни объектов;
    \item средства работы с памятью;
    \item средства обработки ошибок и исключительных ситуаций;
    \item механизмы подключения библиотек.
\end{itemize}

К процедурным языкам обычно относят C, Pascal, Ada, Modula-2, Fortran, а также процедурные подмножества многих современных языков.

\subsection{Императивная модель вычислений}

Процедурное программирование является разновидностью \textbf{императивного программирования}. В императивной модели программа задаёт не только то, какой результат нужно получить, но и \textbf{как именно} его вычислить.

Состояние программы можно представить как отображение
\[
    S : V \to D,
\]
где $V$ --- множество переменных, а $D$ --- множество возможных значений.

Команда изменяет состояние:
\[
    C : S \to S'.
\]

Например, присваивание
\[
    x := x + 1
\]
преобразует состояние $S$ в состояние $S'$, где
\[
    S'(x) = S(x) + 1,
\]
а значения остальных переменных, как правило, остаются неизменными.

\subsection{Основные управляющие конструкции}

В структурном программировании выделяют три базовые управляющие конструкции:
\begin{enumerate}
    \item последовательность;
    \item ветвление;
    \item цикл.
\end{enumerate}

\subsubsection{Последовательность}

\textbf{Последовательность} означает выполнение операторов один за другим в порядке записи.

\[
\begin{array}{l}
    S_1;\\
    S_2;\\
    \dots\\
    S_n.
\end{array}
\]

Если оператор $S_1$ переводит состояние $\sigma_0$ в $\sigma_1$, оператор $S_2$ --- $\sigma_1$ в $\sigma_2$, и так далее, то вся последовательность переводит
\[
    \sigma_0 \to \sigma_n.
\]

\paragraph{Пример.}

\begin{verbatim}
x := 2;
y := x + 3;
z := y * 2;
\end{verbatim}

После выполнения:
\[
    x = 2,\quad y = 5,\quad z = 10.
\]

\subsubsection{Ветвление}

\textbf{Ветвление} выбирает одну из альтернатив в зависимости от значения логического условия.

Общая форма:

\begin{verbatim}
if condition then
    S1
else
    S2
end if;
\end{verbatim}

Семантика:

\[
\text{if } B \text{ then } S_1 \text{ else } S_2
\]

означает:
\[
\begin{cases}
S_1, & \text{если } B = \text{true},\\
S_2, & \text{если } B = \text{false}.
\end{cases}
\]

\paragraph{Пример.}

\begin{verbatim}
if x >= 0 then
    abs_x := x;
else
    abs_x := -x;
end if;
\end{verbatim}

\subsubsection{Множественный выбор}

Для выбора одной из нескольких ветвей часто используется конструкция \texttt{case} или \texttt{switch}.

\begin{verbatim}
case day of
    1: name := "Monday";
    2: name := "Tuesday";
    3: name := "Wednesday";
    ...
end case;
\end{verbatim}

Такая конструкция удобна, когда выбор зависит от значения одной переменной дискретного типа.

\subsubsection{Циклы}

\textbf{Цикл} предназначен для многократного выполнения фрагмента программы.

Основные виды циклов:
\begin{itemize}
    \item цикл с предусловием;
    \item цикл с постусловием;
    \item цикл со счётчиком.
\end{itemize}

\paragraph{Цикл с предусловием.}

\begin{verbatim}
while condition do
    S
end while;
\end{verbatim}

Оператор $S$ выполняется, пока условие истинно. Если условие изначально ложно, тело цикла не выполнится ни разу.

\paragraph{Пример.}

\begin{verbatim}
i := 1;
sum := 0;
while i <= n do
    sum := sum + i;
    i := i + 1;
end while;
\end{verbatim}

После выполнения:
\[
    sum = 1 + 2 + \dots + n = \frac{n(n+1)}{2}.
\]

\paragraph{Цикл с постусловием.}

\begin{verbatim}
repeat
    S
until condition;
\end{verbatim}

Тело цикла выполняется хотя бы один раз, затем проверяется условие выхода.

\paragraph{Пример.}

\begin{verbatim}
repeat
    read(x);
until x >= 0;
\end{verbatim}

Цикл повторяет ввод, пока не будет введено неотрицательное число.

\paragraph{Цикл со счётчиком.}

\begin{verbatim}
for i := 1 to n do
    S
end for;
\end{verbatim}

Используется, когда заранее известно количество повторений.

\paragraph{Пример.}

\begin{verbatim}
factorial := 1;
for i := 1 to n do
    factorial := factorial * i;
end for;
\end{verbatim}

После выполнения:
\[
    factorial = n!.
\]

\subsection{Теорема Бёма--Якопини}

\textbf{Теорема Бёма--Якопини.}
Всякая вычислимая функция, реализуемая блок-схемой с произвольными переходами, может быть реализована программой, использующей только три управляющие конструкции:
\begin{enumerate}
    \item последовательность;
    \item ветвление;
    \item цикл.
\end{enumerate}

\subsubsection{Идея доказательства}

Пусть дана произвольная блок-схема, состоящая из конечного числа блоков:
\[
    B_1, B_2, \dots, B_n.
\]

Каждому блоку поставим в соответствие номер. Введём специальную переменную управления:
\[
    pc,
\]
где $pc$ --- номер текущего выполняемого блока. Такая переменная называется \textbf{счётчиком команд}.

Вместо произвольных переходов построим один цикл:

\begin{verbatim}
pc := 1;
while pc != 0 do
    case pc of
        1: выполнить блок B1; изменить pc;
        2: выполнить блок B2; изменить pc;
        ...
        n: выполнить блок Bn; изменить pc;
    end case;
end while;
\end{verbatim}

Значение $pc = 0$ означает завершение программы.

Каждый переход из исходной блок-схемы заменяется присваиванием нового значения переменной $pc$.

Например, если после блока $B_i$ нужно перейти к $B_j$, то записывается:
\[
    pc := j.
\]

Если блок является условным, например:

\[
    \text{если } x > 0 \text{ перейти к } B_j, \text{ иначе к } B_k,
\]

то он заменяется на:

\begin{verbatim}
if x > 0 then
    pc := j;
else
    pc := k;
end if;
\end{verbatim}

Таким образом, вся исходная программа с произвольными переходами моделируется с помощью:
\begin{itemize}
    \item одного цикла \texttt{while};
    \item множественного выбора \texttt{case};
    \item условного оператора \texttt{if};
    \item последовательного выполнения операторов.
\end{itemize}

Следовательно, произвольные операторы перехода не являются принципиально необходимыми для выражения алгоритмов. Теорема обосновывает возможность структурного программирования.

\subsection{Структура программы}

Структура процедурной программы зависит от языка, но обычно включает следующие части:

\begin{enumerate}
    \item заголовок программы или модуля;
    \item подключение библиотек;
    \item объявления констант;
    \item объявления типов;
    \item объявления глобальных переменных;
    \item объявления процедур и функций;
    \item основную часть программы.
\end{enumerate}

\subsubsection{Обобщённая структура программы}

\begin{verbatim}
program ProgramName;

import Libraries;

const
    declarations_of_constants;

type
    declarations_of_types;

var
    declarations_of_variables;

procedure P(...);
begin
    ...
end;

function F(...): ReturnType;
begin
    ...
end;

begin
    main_program_body;
end.
\end{verbatim}

\subsubsection{Пример на псевдо-Pascal}

\begin{verbatim}
program Example;

const
    MaxN = 100;

type
    IntArray = array[1..MaxN] of integer;

var
    a: IntArray;
    n: integer;

function Sum(n: integer; a: IntArray): integer;
var
    i, s: integer;
begin
    s := 0;
    for i := 1 to n do
        s := s + a[i];
    Sum := s;
end;

begin
    read(n);
    ...
    write(Sum(n, a));
end.
\end{verbatim}

\subsection{Работа с данными}

\subsubsection{Переменные}

\textbf{Переменная} --- именованная область памяти, в которой хранится значение некоторого типа.

Переменная характеризуется:
\begin{itemize}
    \item именем;
    \item типом;
    \item значением;
    \item адресом;
    \item временем жизни;
    \item областью видимости.
\end{itemize}

Объявление переменной связывает имя с типом:

\begin{verbatim}
var
    x: integer;
\end{verbatim}

В языках семейства C:

\begin{verbatim}
int x;
\end{verbatim}

\subsubsection{Присваивание}

Оператор присваивания изменяет значение переменной.

\begin{verbatim}
x := x + 1;
\end{verbatim}

Сначала вычисляется выражение справа, затем результат помещается в переменную слева.

Если до выполнения:
\[
    x = 5,
\]
то после:
\[
    x = 6.
\]

\subsubsection{Константы}

\textbf{Константа} --- именованное значение, которое не изменяется во время выполнения программы.

\begin{verbatim}
const
    Pi = 3.1415926;
    MaxSize = 1000;
\end{verbatim}

Использование констант улучшает читаемость и сопровождаемость программы.

\subsection{Типы данных}

\textbf{Тип данных} определяет:
\begin{itemize}
    \item множество допустимых значений;
    \item множество допустимых операций;
    \item способ представления значений в памяти.
\end{itemize}

Например, тип \texttt{integer} задаёт целые числа и операции сложения, вычитания, умножения, сравнения и т.д.

\subsubsection{Булевский тип}

\textbf{Булевский тип} имеет два значения:
\[
    \text{true}, \quad \text{false}.
\]

Основные операции:
\begin{itemize}
    \item отрицание: $\neg a$;
    \item конъюнкция: $a \land b$;
    \item дизъюнкция: $a \lor b$;
    \item исключающее или: $a \oplus b$.
\end{itemize}

\paragraph{Таблица истинности для конъюнкции.}

\[
\begin{array}{c|c|c}
a & b & a \land b\\
\hline
false & false & false\\
false & true & false\\
true & false & false\\
true & true & true
\end{array}
\]

\paragraph{Пример.}

\begin{verbatim}
if age >= 18 and has_passport then
    allow := true;
else
    allow := false;
end if;
\end{verbatim}

\subsubsection{Целочисленные типы}

\textbf{Целочисленные типы} предназначены для хранения целых чисел.

Типичный набор:
\begin{itemize}
    \item \texttt{short};
    \item \texttt{integer};
    \item \texttt{long};
    \item \texttt{unsigned integer}.
\end{itemize}

Если для хранения целого числа используется $k$ бит, то возможны:
\[
    2^k
\]
различных битовых комбинаций.

Для беззнакового представления диапазон:
\[
    0 \dots 2^k - 1.
\]

Для знакового представления в дополнительном коде:
\[
    -2^{k-1} \dots 2^{k-1} - 1.
\]

\paragraph{Пример для 8 бит.}

Беззнаковое число:
\[
    0 \dots 255.
\]

Знаковое число:
\[
    -128 \dots 127.
\]

\subsubsection{Переполнение}

\textbf{Переполнение} возникает, когда результат операции не помещается в диапазон типа.

Например, если переменная типа \texttt{unsigned byte} имеет диапазон $0 \dots 255$, то операция:
\[
    255 + 1
\]
даёт математический результат $256$, который не представим в этом типе.

В разных языках переполнение может:
\begin{itemize}
    \item приводить к ошибке;
    \item давать результат по модулю $2^k$;
    \item считаться неопределённым поведением.
\end{itemize}

\subsubsection{Типы с плавающей точкой}

\textbf{Типы с плавающей точкой} используются для приближённого представления вещественных чисел.

Обычно число представляется в форме:
\[
    x = (-1)^s \cdot m \cdot b^e,
\]
где:
\begin{itemize}
    \item $s$ --- бит знака;
    \item $m$ --- мантисса;
    \item $b$ --- основание системы счисления, обычно $2$;
    \item $e$ --- порядок.
\end{itemize}

Например:
\[
    12.5 = 1.25 \cdot 10^1.
\]

В двоичной системе:
\[
    12.5_{10} = 1100.1_2 = 1.1001_2 \cdot 2^3.
\]

\paragraph{Особенности чисел с плавающей точкой.}

\begin{itemize}
    \item Не все вещественные числа представимы точно.
    \item Операции могут давать округлённый результат.
    \item Сравнение на равенство может быть ненадёжным.
    \item Возможны специальные значения: бесконечность, NaN.
\end{itemize}

\paragraph{Пример проблемы точности.}

Число $0.1$ не имеет конечного двоичного представления, поэтому выражение:

\begin{verbatim}
0.1 + 0.2 = 0.3
\end{verbatim}

в машине может оказаться ложным из-за округления.

Для сравнения часто используют допуск $\varepsilon$:

\[
    |a - b| < \varepsilon.
\]

\subsubsection{Символьный тип}

\textbf{Символьный тип} предназначен для хранения отдельных символов.

Примеры значений:
\[
    'a', \quad 'Z', \quad '7', \quad '+', \quad '\textbackslash n'.
\]

Символы кодируются числами. Например, в ASCII символ \texttt{'A'} имеет код $65$, а символ \texttt{'a'} --- код $97$.

\paragraph{Пример.}

\begin{verbatim}
ch := 'A';
if ch >= 'A' and ch <= 'Z' then
    is_upper := true;
end if;
\end{verbatim}

\subsubsection{Типы диапазона}

\textbf{Тип диапазона} --- тип, множество значений которого является подмножеством значений другого дискретного типа.

Например:

\begin{verbatim}
type Day = 1..31;
type Month = 1..12;
\end{verbatim}

Значения типа \texttt{Month} могут принимать только значения:
\[
    1,2,\dots,12.
\]

Типы диапазона полезны, поскольку:
\begin{itemize}
    \item повышают выразительность программы;
    \item позволяют компилятору или среде выполнения обнаруживать ошибки;
    \item документируют ограничения предметной области.
\end{itemize}

\paragraph{Пример.}

\begin{verbatim}
var
    m: Month;

m := 5;   { корректно }
m := 15;  { ошибка диапазона }
\end{verbatim}

\subsubsection{Перечислимые типы}

\textbf{Перечислимый тип} задаётся явным перечислением всех возможных значений.

\begin{verbatim}
type Color = (Red, Green, Blue);
type WeekDay = (Monday, Tuesday, Wednesday,
                Thursday, Friday, Saturday, Sunday);
\end{verbatim}

Перечислимые типы улучшают читаемость программы.

\paragraph{Пример.}

\begin{verbatim}
var
    c: Color;

c := Red;

case c of
    Red:   write("red");
    Green: write("green");
    Blue:  write("blue");
end case;
\end{verbatim}

\subsubsection{Указатели}

\textbf{Указатель} --- значение, содержащее адрес объекта в памяти.

Если переменная $x$ расположена по адресу $a$, то указатель $p$ может хранить этот адрес:
\[
    p = \&x.
\]

Операции с указателями:
\begin{itemize}
    \item взятие адреса;
    \item разыменование;
    \item сравнение указателей;
    \item присваивание указателей;
    \item в некоторых языках --- адресная арифметика.
\end{itemize}

\paragraph{Пример на C-подобном языке.}

\begin{verbatim}
int x = 10;
int *p = &x;

*p = 20;
\end{verbatim}

После выполнения:
\[
    x = 20.
\]

\paragraph{Схема.}

\[
\begin{array}{c|c}
\text{имя} & \text{содержимое}\\
\hline
x & 10\\
p & \text{адрес } x
\end{array}
\]

После операции \texttt{*p = 20}:

\[
\begin{array}{c|c}
\text{имя} & \text{содержимое}\\
\hline
x & 20\\
p & \text{адрес } x
\end{array}
\]

\paragraph{Нулевой указатель.}

\textbf{Нулевой указатель} не указывает ни на какой объект.

\begin{verbatim}
p := null;
\end{verbatim}

Попытка разыменовать нулевой указатель обычно является ошибкой.

\paragraph{Опасности указателей.}

\begin{itemize}
    \item разыменование нулевого указателя;
    \item висячие указатели;
    \item утечки памяти;
    \item выход за границы массива;
    \item двойное освобождение памяти.
\end{itemize}

\subsection{Структуры данных}

\subsubsection{Массивы}

\textbf{Массив} --- структура данных, состоящая из элементов одного типа, доступ к которым осуществляется по индексу.

Одномерный массив:
\[
    A[1], A[2], \dots, A[n].
\]

Объявление:

\begin{verbatim}
var
    a: array[1..10] of integer;
\end{verbatim}

В C-подобном языке:

\begin{verbatim}
int a[10];
\end{verbatim}

\paragraph{Индексация.}

В разных языках индексация может начинаться:
\begin{itemize}
    \item с $0$;
    \item с $1$;
    \item с произвольной нижней границы.
\end{itemize}

\paragraph{Пример.}

\begin{verbatim}
a[1] := 10;
a[2] := 20;
a[3] := a[1] + a[2];
\end{verbatim}

После выполнения:
\[
    a[3] = 30.
\]

\subsubsection{Многомерные массивы}

\textbf{Двумерный массив} можно рассматривать как таблицу:

\[
    A =
    \begin{pmatrix}
    a_{11} & a_{12} & a_{13}\\
    a_{21} & a_{22} & a_{23}
    \end{pmatrix}.
\]

Объявление:

\begin{verbatim}
var
    matrix: array[1..2, 1..3] of integer;
\end{verbatim}

Обращение к элементу:
\[
    matrix[i,j].
\]

\subsubsection{Алгоритм суммирования элементов массива}

\paragraph{Задача.}
Дан массив $A[1..n]$. Найти сумму его элементов.

\paragraph{Идея.}
Завести переменную $s$, сначала равную $0$, и последовательно прибавлять к ней элементы массива.

\paragraph{Шаги алгоритма.}

\begin{enumerate}
    \item Присвоить $s := 0$.
    \item Присвоить $i := 1$.
    \item Если $i > n$, перейти к шагу 6.
    \item Выполнить $s := s + A[i]$.
    \item Выполнить $i := i + 1$ и перейти к шагу 3.
    \item Выдать $s$ как результат.
\end{enumerate}

\paragraph{Псевдокод.}

\begin{verbatim}
s := 0;
for i := 1 to n do
    s := s + A[i];
return s;
\end{verbatim}

\paragraph{Пример.}

Пусть:
\[
    A = [3, 5, 2, 4].
\]

Тогда:

\[
\begin{array}{c|c|c}
i & A[i] & s\\
\hline
0 & - & 0\\
1 & 3 & 3\\
2 & 5 & 8\\
3 & 2 & 10\\
4 & 4 & 14
\end{array}
\]

Ответ:
\[
    s = 14.
\]

\paragraph{Сложность.}
Алгоритм выполняет $n$ сложений, поэтому временная сложность:
\[
    O(n).
\]

Дополнительная память:
\[
    O(1).
\]

\subsubsection{Записи}

\textbf{Запись} --- структура данных, объединяющая несколько полей, возможно разных типов.

В Pascal-подобном языке:

\begin{verbatim}
type Student = record
    name: string;
    age: integer;
    average_grade: real;
end;
\end{verbatim}

В C-подобном языке аналогом является \texttt{struct}:

\begin{verbatim}
struct Student {
    char name[100];
    int age;
    double average_grade;
};
\end{verbatim}

\paragraph{Обращение к полям.}

\begin{verbatim}
s.name := "Ivan";
s.age := 20;
s.average_grade := 4.7;
\end{verbatim}

\paragraph{Схема записи.}

\[
\begin{array}{c|c}
\text{поле} & \text{значение}\\
\hline
name & \text{Ivan}\\
age & 20\\
average\_grade & 4.7
\end{array}
\]

\subsubsection{Массивы записей}

Часто записи объединяются в массивы.

\begin{verbatim}
var
    group: array[1..30] of Student;
\end{verbatim}

Тогда:
\[
    group[i].name
\]
обозначает имя $i$-го студента.

\subsection{Процедуры и функции}

\subsubsection{Понятие подпрограммы}

\textbf{Подпрограмма} --- именованный фрагмент программы, который можно вызвать из других частей программы.

Различают:
\begin{itemize}
    \item \textbf{процедуры} --- выполняют действия, но не обязательно возвращают значение;
    \item \textbf{функции} --- вычисляют и возвращают значение.
\end{itemize}

\subsubsection{Процедуры}

Процедура имеет имя, параметры и тело.

\begin{verbatim}
procedure PrintHello();
begin
    write("Hello");
end;
\end{verbatim}

Вызов:

\begin{verbatim}
PrintHello();
\end{verbatim}

\subsubsection{Функции}

Функция возвращает значение.

\begin{verbatim}
function Square(x: integer): integer;
begin
    Square := x * x;
end;
\end{verbatim}

Вызов:

\begin{verbatim}
y := Square(5);
\end{verbatim}

После выполнения:
\[
    y = 25.
\]

\subsubsection{Стек вызовов}

При вызове процедуры обычно создаётся \textbf{активационная запись} в стеке вызовов.

Она может содержать:
\begin{itemize}
    \item параметры;
    \item локальные переменные;
    \item адрес возврата;
    \item сохранённые регистры;
    \item служебную информацию.
\end{itemize}

Если процедура $A$ вызывает $B$, а $B$ вызывает $C$, стек можно представить так:

\[
\begin{array}{|c|}
\hline
\text{активационная запись } C\\
\hline
\text{активационная запись } B\\
\hline
\text{активационная запись } A\\
\hline
\text{основная программа}\\
\hline
\end{array}
\]

После завершения $C$ её запись удаляется из стека, управление возвращается в $B$.

\subsection{Передача параметров}

\textbf{Фактические параметры} --- аргументы, указанные при вызове.

\textbf{Формальные параметры} --- параметры, указанные в объявлении процедуры или функции.

Например:

\begin{verbatim}
procedure P(x: integer);
begin
    ...
end;

P(10);
\end{verbatim}

Здесь \texttt{x} --- формальный параметр, а \texttt{10} --- фактический параметр.

\subsubsection{Передача по значению}

При \textbf{передаче по значению} формальный параметр получает копию значения фактического параметра.

Изменение формального параметра не изменяет фактический параметр.

\paragraph{Пример.}

\begin{verbatim}
procedure Inc(x: integer);
begin
    x := x + 1;
end;

a := 5;
Inc(a);
\end{verbatim}

После вызова:
\[
    a = 5.
\]

Изменилась только локальная копия \texttt{x}.

\paragraph{Схема.}

До вызова:

\[
\begin{array}{c|c}
a & 5
\end{array}
\]

Во время вызова:

\[
\begin{array}{c|c}
a & 5\\
x & 5
\end{array}
\]

После \texttt{x := x + 1}:

\[
\begin{array}{c|c}
a & 5\\
x & 6
\end{array}
\]

После выхода из процедуры \texttt{x} уничтожается, а \texttt{a} остаётся равной $5$.

\subsubsection{Передача по ссылке}

При \textbf{передаче по ссылке} формальный параметр становится другим именем для фактического параметра. Изменение формального параметра изменяет фактический параметр.

\paragraph{Пример.}

\begin{verbatim}
procedure Inc(var x: integer);
begin
    x := x + 1;
end;

a := 5;
Inc(a);
\end{verbatim}

После вызова:
\[
    a = 6.
\]

\paragraph{Схема.}

\[
\begin{array}{c|c}
a, x & 5
\end{array}
\]

После операции:

\[
\begin{array}{c|c}
a, x & 6
\end{array}
\]

Имена \texttt{a} и \texttt{x} обозначают один и тот же объект.

\subsubsection{Передача по результату}

При \textbf{передаче по результату} значение фактического параметра при входе в процедуру не используется. Формальный параметр работает как локальная переменная, а при выходе его значение копируется в фактический параметр.

Такой механизм иногда называют \textbf{copy-out}.

\paragraph{Пример.}

\begin{verbatim}
procedure SetToTen(out x: integer);
begin
    x := 10;
end;

a := 5;
SetToTen(a);
\end{verbatim}

После вызова:
\[
    a = 10.
\]

\paragraph{Схема.}

До вызова:

\[
\begin{array}{c|c}
a & 5
\end{array}
\]

Во время выполнения:

\[
\begin{array}{c|c}
a & 5\\
x & \text{неопределено}
\end{array}
\]

После присваивания:

\[
\begin{array}{c|c}
a & 5\\
x & 10
\end{array}
\]

При выходе из процедуры:

\[
\begin{array}{c|c}
a & 10
\end{array}
\]

\subsubsection{Передача по значению-результату}

Близкий механизм --- \textbf{передача по значению-результату}, или \textbf{copy-in/copy-out}.

При входе значение фактического параметра копируется в формальный параметр. При выходе значение формального параметра копируется обратно.

\paragraph{Пример.}

\begin{verbatim}
procedure AddOne(inout x: integer);
begin
    x := x + 1;
end;

a := 5;
AddOne(a);
\end{verbatim}

После вызова:
\[
    a = 6.
\]

\subsubsection{Сравнение способов передачи параметров}

\[
\begin{array}{c|c|c|c}
\text{Способ} & \text{Входное значение} & \text{Изменяет аргумент} & \text{Типичный смысл}\\
\hline
\text{по значению} & \text{копируется} & \text{нет} & \text{входной параметр}\\
\text{по ссылке} & \text{используется объект} & \text{да} & \text{модификация объекта}\\
\text{по результату} & \text{не используется} & \text{да} & \text{выходной параметр}\\
\text{значение-результат} & \text{копируется} & \text{да} & \text{вход-выход}
\end{array}
\]

\subsubsection{Проблема наложения параметров}

Если один и тот же фактический параметр передаётся нескольким формальным параметрам, может возникнуть неоднозначность.

\paragraph{Пример.}

\begin{verbatim}
procedure P(var x: integer; var y: integer);
begin
    x := x + 1;
    y := y + 1;
end;

a := 5;
P(a, a);
\end{verbatim}

Если оба параметра являются ссылками на \texttt{a}, то:
\[
    a = 7.
\]

Но при механизме копирования результата порядок обратного копирования может влиять на результат.

\subsection{Локализация переменных}

\subsubsection{Область видимости}

\textbf{Область видимости} имени --- часть программы, в которой это имя доступно.

Переменные бывают:
\begin{itemize}
    \item глобальные;
    \item локальные;
    \item статические локальные;
    \item динамические.
\end{itemize}

\subsubsection{Глобальные переменные}

\textbf{Глобальная переменная} объявляется вне процедур и доступна во многих частях программы.

\begin{verbatim}
var
    counter: integer;

procedure Inc();
begin
    counter := counter + 1;
end;
\end{verbatim}

Преимущество:
\begin{itemize}
    \item удобство доступа.
\end{itemize}

Недостатки:
\begin{itemize}
    \item труднее понимать зависимости;
    \item сложнее тестировать;
    \item возможны скрытые побочные эффекты.
\end{itemize}

\subsubsection{Локальные переменные}

\textbf{Локальная переменная} объявляется внутри процедуры или блока и доступна только внутри него.

\begin{verbatim}
procedure P();
var
    x: integer;
begin
    x := 10;
end;
\end{verbatim}

После завершения процедуры локальная автоматическая переменная обычно уничтожается.

\subsubsection{Сокрытие имён}

Если локальная переменная имеет то же имя, что и глобальная, локальная переменная может скрывать глобальную.

\begin{verbatim}
var
    x: integer;

procedure P();
var
    x: integer;
begin
    x := 10;
end;
\end{verbatim}

Внутри \texttt{P} имя \texttt{x} относится к локальной переменной.

\subsubsection{Статическая и динамическая области видимости}

При \textbf{статической}, или \textbf{лексической}, области видимости связь имени с объявлением определяется текстом программы.

При \textbf{динамической} области видимости связь имени определяется цепочкой вызовов во время выполнения.

Большинство современных процедурных языков используют статическую область видимости.

\subsection{Побочные эффекты}

\textbf{Побочный эффект} --- изменение состояния, не выраженное явно в возвращаемом значении функции.

Примеры побочных эффектов:
\begin{itemize}
    \item изменение глобальной переменной;
    \item изменение параметра, переданного по ссылке;
    \item ввод-вывод;
    \item изменение файла;
    \item выделение или освобождение памяти;
    \item генерация исключения.
\end{itemize}

\paragraph{Пример.}

\begin{verbatim}
var
    counter: integer;

function Next(): integer;
begin
    counter := counter + 1;
    Next := counter;
end;
\end{verbatim}

Функция \texttt{Next} не только возвращает значение, но и изменяет глобальную переменную \texttt{counter}.

\subsubsection{Чистые функции}

\textbf{Чистая функция} --- функция, которая:
\begin{itemize}
    \item при одинаковых аргументах всегда возвращает одинаковый результат;
    \item не имеет побочных эффектов.
\end{itemize}

Пример:

\begin{verbatim}
function Square(x: integer): integer;
begin
    Square := x * x;
end;
\end{verbatim}

Чистые функции проще тестировать, оптимизировать и использовать в доказательствах корректности программ.

\subsection{Рекурсия}

\textbf{Рекурсия} --- ситуация, когда процедура или функция вызывает саму себя.

\subsubsection{Факториал}

Математическое определение:
\[
    n! =
    \begin{cases}
        1, & n = 0,\\
        n \cdot (n-1)!, & n > 0.
    \end{cases}
\]

\paragraph{Рекурсивная функция.}

\begin{verbatim}
function Fact(n: integer): integer;
begin
    if n = 0 then
        Fact := 1;
    else
        Fact := n * Fact(n - 1);
end;
\end{verbatim}

\paragraph{Вычисление $Fact(4)$.}

\[
\begin{aligned}
Fact(4) &= 4 \cdot Fact(3)\\
        &= 4 \cdot 3 \cdot Fact(2)\\
        &= 4 \cdot 3 \cdot 2 \cdot Fact(1)\\
        &= 4 \cdot 3 \cdot 2 \cdot 1 \cdot Fact(0)\\
        &= 4 \cdot 3 \cdot 2 \cdot 1 \cdot 1\\
        &= 24.
\end{aligned}
\]

\subsubsection{Условия корректной рекурсии}

Рекурсивная процедура должна иметь:
\begin{itemize}
    \item базовый случай;
    \item рекурсивный переход;
    \item продвижение к базовому случаю.
\end{itemize}

Если базовый случай отсутствует или недостижим, возникает бесконечная рекурсия и переполнение стека.

\subsection{Обработка исключительных ситуаций}

\subsubsection{Понятие исключительной ситуации}

\textbf{Исключительная ситуация} --- необычное или ошибочное состояние, возникающее при выполнении программы и нарушающее нормальный ход вычислений.

Примеры:
\begin{itemize}
    \item деление на ноль;
    \item выход за границы массива;
    \item ошибка выделения памяти;
    \item ошибка открытия файла;
    \item неверный формат входных данных;
    \item разыменование нулевого указателя.
\end{itemize}

\subsubsection{Подходы к обработке ошибок}

Основные способы обработки ошибок:
\begin{enumerate}
    \item возврат специального кода ошибки;
    \item использование глобальной переменной ошибки;
    \item передача ошибки через выходной параметр;
    \item механизм исключений;
    \item аварийное завершение программы.
\end{enumerate}

\subsubsection{Коды ошибок}

Процедура может возвращать код состояния.

\begin{verbatim}
result := OpenFile("data.txt");
if result != OK then
    write("Error");
end if;
\end{verbatim}

Недостаток: программист может забыть проверить код ошибки.

\subsubsection{Исключения}

\textbf{Исключение} --- объект или событие, сигнализирующее об ошибке и передающее управление специальному обработчику.

Типичная конструкция:

\begin{verbatim}
try
    risky_operation();
catch SomeException
    handle_error();
end try;
\end{verbatim}

В Ada-подобной форме:

\begin{verbatim}
begin
    X := A / B;
exception
    when Constraint_Error =>
        Put_Line("Division error");
end;
\end{verbatim}

\subsubsection{Механизм распространения исключений}

Если исключение возникло внутри процедуры и не обработано там, оно передаётся вызывающей процедуре. Этот процесс продолжается вверх по стеку вызовов.

Пусть:
\[
    A \to B \to C,
\]
где $A$ вызывает $B$, а $B$ вызывает $C$.

Если в $C$ возникает исключение:
\begin{itemize}
    \item сначала ищется обработчик в $C$;
    \item если его нет, управление передаётся в $B$;
    \item если его нет в $B$, управление передаётся в $A$;
    \item если обработчика нет нигде, программа аварийно завершается.
\end{itemize}

\paragraph{Схема.}

\[
\begin{array}{|c|}
\hline
C \quad \text{исключение}\\
\hline
B \quad \text{нет обработчика}\\
\hline
A \quad \text{обработчик найден}\\
\hline
\end{array}
\]

\subsubsection{Алгоритм поиска обработчика исключения}

\paragraph{Шаги.}

\begin{enumerate}
    \item Возникает исключение $E$ в текущей активационной записи.
    \item Проверяется, есть ли в текущем блоке обработчик для $E$.
    \item Если обработчик найден, управление передаётся ему.
    \item Если обработчик не найден, текущая активационная запись удаляется из стека.
    \item Поиск продолжается в вызывающем блоке.
    \item Если стек вызовов исчерпан, программа завершается с необработанным исключением.
\end{enumerate}

\paragraph{Минимальный пример.}

\begin{verbatim}
procedure C();
begin
    raise Error;
end;

procedure B();
begin
    C();
end;

procedure A();
begin
    B();
exception
    when Error =>
        write("handled");
end;
\end{verbatim}

Здесь исключение возникает в \texttt{C}, проходит через \texttt{B} и обрабатывается в \texttt{A}.

\subsubsection{Преимущества исключений}

\begin{itemize}
    \item отделяют основной алгоритм от обработки ошибок;
    \item позволяют передавать ошибку через несколько уровней вызовов;
    \item уменьшают количество проверок после каждого вызова;
    \item делают ошибки явными.
\end{itemize}

\subsubsection{Недостатки исключений}

\begin{itemize}
    \item усложняют поток управления;
    \item могут скрывать места возникновения ошибки;
    \item требуют аккуратного освобождения ресурсов;
    \item могут иметь значительную стоимость при возникновении.
\end{itemize}

\subsubsection{Освобождение ресурсов}

При исключениях важно гарантировать освобождение ресурсов:
\begin{itemize}
    \item файлов;
    \item блокировок;
    \item динамической памяти;
    \item сетевых соединений.
\end{itemize}

Типичный подход:

\begin{verbatim}
try
    acquire_resource();
    use_resource();
finally
    release_resource();
end try;
\end{verbatim}

Блок \texttt{finally} выполняется независимо от того, произошло исключение или нет.

\subsection{Библиотеки процедур}

\subsubsection{Понятие библиотеки}

\textbf{Библиотека} --- набор готовых процедур, функций, типов и констант, предназначенных для повторного использования.

Библиотеки позволяют:
\begin{itemize}
    \item не писать заново стандартные алгоритмы;
    \item использовать проверенный код;
    \item разделять программу на модули;
    \item ускорять разработку.
\end{itemize}

\subsubsection{Виды библиотек}

\begin{enumerate}
    \item \textbf{Стандартные библиотеки} --- входят в состав языка или его реализации.
    \item \textbf{Системные библиотеки} --- предоставляются операционной системой.
    \item \textbf{Пользовательские библиотеки} --- создаются разработчиком.
    \item \textbf{Сторонние библиотеки} --- распространяются отдельно.
\end{enumerate}

\subsubsection{Примеры библиотечных возможностей}

Стандартные библиотеки могут содержать:
\begin{itemize}
    \item ввод-вывод;
    \item математические функции;
    \item работу со строками;
    \item работу с файлами;
    \item управление памятью;
    \item обработку дат и времени;
    \item генерацию случайных чисел.
\end{itemize}

\subsubsection{Подключение библиотек}

В C:

\begin{verbatim}
#include <stdio.h>
#include <math.h>
\end{verbatim}

В Pascal:

\begin{verbatim}
uses Math, SysUtils;
\end{verbatim}

В Ada:

\begin{verbatim}
with Ada.Text_IO;
use Ada.Text_IO;
\end{verbatim}

\subsubsection{Раздельная компиляция}

Большие программы обычно разбиваются на несколько файлов или модулей.

Типичный процесс:
\begin{enumerate}
    \item каждый модуль компилируется отдельно;
    \item компилятор создаёт объектные файлы;
    \item компоновщик объединяет объектные файлы и библиотеки;
    \item получается исполняемая программа.
\end{enumerate}

\[
\text{исходные файлы}
\to
\text{объектные файлы}
\to
\text{исполняемый файл}.
\]

\subsubsection{Интерфейс и реализация}

Хорошая библиотека разделяет:
\begin{itemize}
    \item \textbf{интерфейс} --- что библиотека предоставляет;
    \item \textbf{реализацию} --- как это устроено внутри.
\end{itemize}

Например, интерфейс может содержать объявление функции:

\begin{verbatim}
function Max(a, b: integer): integer;
\end{verbatim}

А реализация:

\begin{verbatim}
function Max(a, b: integer): integer;
begin
    if a > b then
        Max := a;
    else
        Max := b;
end;
\end{verbatim}

Пользователь библиотеки знает, как вызвать функцию, но не обязан знать детали её реализации.

\subsection{Минимальный пример процедурной программы}

\subsubsection{Задача}

Дан массив целых чисел. Нужно найти максимальный элемент.

\subsubsection{Алгоритм}

\begin{enumerate}
    \item Считать $n$.
    \item Считать массив $A[1..n]$.
    \item Присвоить $max := A[1]$.
    \item Для каждого $i$ от $2$ до $n$:
    \begin{enumerate}
        \item если $A[i] > max$, то присвоить $max := A[i]$.
    \end{enumerate}
    \item Вывести $max$.
\end{enumerate}

\subsubsection{Псевдокод}

\begin{verbatim}
function MaxElement(A: array[1..n] of integer): integer;
var
    i, max: integer;
begin
    max := A[1];

    for i := 2 to n do
        if A[i] > max then
            max := A[i];

    MaxElement := max;
end;
\end{verbatim}

\subsubsection{Пример работы}

Пусть:
\[
    A = [4, 1, 9, 3, 7].
\]

\[
\begin{array}{c|c|c}
i & A[i] & max\\
\hline
1 & 4 & 4\\
2 & 1 & 4\\
3 & 9 & 9\\
4 & 3 & 9\\
5 & 7 & 9
\end{array}
\]

Ответ:
\[
    9.
\]

\subsubsection{Корректность алгоритма}

Докажем корректность алгоритма поиска максимума.

\paragraph{Инвариант цикла.}

После обработки элементов
\[
    A[1], A[2], \dots, A[i]
\]
переменная \texttt{max} содержит максимум среди этих элементов:
\[
    max = \max\{A[1], A[2], \dots, A[i]\}.
\]

\paragraph{База.}

Перед началом цикла:
\[
    max := A[1].
\]

Следовательно,
\[
    max = \max\{A[1]\}.
\]

Инвариант верен для $i = 1$.

\paragraph{Переход.}

Пусть после обработки первых $i-1$ элементов:
\[
    max = \max\{A[1], \dots, A[i-1]\}.
\]

На шаге $i$ сравниваются $A[i]$ и \texttt{max}.

Если:
\[
    A[i] > max,
\]
то новое значение:
\[
    max := A[i].
\]

Тогда:
\[
    max = \max\{A[1], \dots, A[i]\}.
\]

Если:
\[
    A[i] \le max,
\]
то \texttt{max} не изменяется, и оно уже не меньше всех элементов $A[1], \dots, A[i-1]$, а также не меньше $A[i]$. Поэтому:
\[
    max = \max\{A[1], \dots, A[i]\}.
\]

Следовательно, инвариант сохраняется.

\paragraph{Завершение.}

После завершения цикла обработаны все элементы:
\[
    A[1], A[2], \dots, A[n].
\]

По инварианту:
\[
    max = \max\{A[1], A[2], \dots, A[n]\}.
\]

Значит, алгоритм возвращает максимальный элемент массива.

\subsubsection{Сложность}

Цикл выполняет $n-1$ сравнений, поэтому временная сложность:
\[
    O(n).
\]

Дополнительная память:
\[
    O(1).
\]

\subsection{Основные достоинства процедурного программирования}

\begin{itemize}
    \item Простая и естественная модель вычислений.
    \item Возможность пошагового описания алгоритма.
    \item Поддержка декомпозиции через процедуры и функции.
    \item Эффективная реализация на машинах фон Неймана.
    \item Удобство для системного и прикладного программирования.
\end{itemize}

\subsection{Основные недостатки процедурного программирования}

\begin{itemize}
    \item При больших программах трудно управлять глобальным состоянием.
    \item Побочные эффекты усложняют отладку.
    \item Сильная зависимость от порядка выполнения операторов.
    \item Не всегда удобно моделировать сложные предметные области.
    \item Ошибки работы с памятью в языках с указателями могут быть трудно обнаружимыми.
\end{itemize}

\subsection{Краткие выводы}

Процедурное программирование основано на представлении программы как последовательности команд, изменяющих состояние. Его основными средствами являются переменные, присваивания, управляющие конструкции, процедуры и функции. Для организации данных применяются простые типы, массивы, записи и указатели. Процедуры позволяют декомпозировать программу, но требуют аккуратного обращения с параметрами, областями видимости и побочными эффектами. Исключения и библиотеки расширяют выразительные возможности языка и позволяют строить более надёжные и модульные программы.

\end{document}
